\cotefile{79}
\scanpage{129}\mspage{65}

\diagram{\begin{tikzcd}[ampersand replacement=\&,column sep=large]
(62)\quad 1 \arrow[r] \& G_0 \arrow[r] \& G \arrow[r,"\delta"] \arrow[d] \& Z' \\
\& \& \mathrm{Cent}(G)=Z' \arrow[ur,swap,"\lambda\mapsto\lambda^2"] \&
\end{tikzcd}}

On voit donc que $\delta(G)$ contient \deleted{$Z'^{2}$}
\[
Z'^{2}=\{\lambda^2\mid\lambda\in Z'\},
\]
mais cela ne signifie pas pour autant que $\delta$ soit surjectif, i.e. que $G$ soit transitif sur l'ens des structures $\Pi$. On voit de suite que l'orbite d'un $i_0\in\unsure{\Pi_S|\Pi_s|\pi}$ sous $G$, correspondant à la structure $A_{i_0}=\unsure{A|A_i|A'|A_0}$ sur $S$, est formée de l'ens des structures $A_\lambda=A_{\lambda\cdot i_0}$ sur $S$ ``équivalentes'' à $A_{i_0}$ (au sens de la p. 61) qui lui sont de plus \uline{isomorphes}. L'axiome (To 5) -- qu'on n'avait pas utilisé jusqu'à maintenant) signifie justement qu'on trouve ainsi tout \deleted{$\unsure{Z}$} $\Pi$, i.e. que (\unsure{63|62|61}) est une suite exacte courte
\[
1\longrightarrow G_0\longrightarrow G\xrightarrow{\delta}Z'\longrightarrow 1\,.\tag*{(63)}
\]

% archive sheet number: 128
\scanpage{130}\mspage{66}

\uline{VIII} \uline{Axiomatique : reformulation}.

Une structure \uline{triangulée} sur un ens $S$ est \inserted{la donnée d'}un ens
\[
\underset{\substack{\text{ens des}\\\text{faces}}}{F_0}\subset\unsure{\mathfrak{P}|\mathcal{P}|P}_3(S)\tag*{(1)}
\]
satisfaisant
\[
S=\bigcup_{f\in F}f\,,\tag*{(2)}
\]
et on lui demande ici de satisfaire les conditions

(3)\qquad (To 1)\qquad $F_0$ est l'ens des faces pour une structure de 2-polyèdre \deleted{\unsure{rég|rég.}} combinatoire \inserted{connexe} ayant $S$ comme ens des sommets, et comme ens des arêtes
\[
A_0=\{a\in\mathfrak{P}_2(S)\mid\exists f\in F_0,\ f\supset a\},
\]
et les relations d'incidences évidentes,

ce qui signifie que pour $\forall\,a\in\mathfrak{P}_2(S)$, s'il existe $f\ni F_0$ avec $f\supset a$, il en existe exactement 2, et que l'on a les conditions de connexité globale (deux \deleted{\unsure{sommets}}\inserted{faces} de $S$ sont reliables par une ``chaîne de faces'') et locale en tt $s\in S$ ($=$ chose, avec les faces $\ni s$). On suppose de plus que $F_0$ est déterminé par $\unsure{A_0|F_0}$ :
\[
\mathrm{To}\,2)\qquad F_0=\{f\in\mathfrak{P}_3(S))\mid\forall a\in\mathfrak{P}_2(f),\ \text{on a }a\in A_0\}.\tag*{(4)}
\]

Considérons maintenant les \deleted{deux}\deleted{\inserted{trois}}\inserted{deux} groupes

% archive sheet number: 129
\scanpage{131}

\[
\deleted{\unsure{Z}}\,G_0=\mathrm{Aut}(S,F_0)=\mathrm{Aut}(S,A_0)=\mathrm{Aut}\,P_0\tag*{(5)}
\]
(où $P_0=(S,A_0,F_0,\text{ relations d'incidences}\ldots)$)
\[
\widehat Z=\mathrm{Centr}_{\unsure{\mathfrak S|\mathcal S}_S}(G_0)\tag*{(6)}
\]
% slanted
\margin{-- \unsure{laisser|je laisse} tomber le (\unsure{2|3|5}) finalement !}

\deleted{\unsure{(7)}} et soit
\[
\begin{aligned}
Z^{\unsure{*|\sharp}}=\bigl\{\lambda\in\widehat Z\ \big|\ &\exists f\in\mathfrak S_S\ \text{t.q. }\forall\{s,t\}\inS,\ \text{on a}\\
&\{s,t\}\in A_0\Longleftrightarrow\{f^{-1}s,f^{-1}\lambda t\}\in A_0\bigr\}.
\end{aligned}\tag*{(7)}
\]

\deleted{i.e.} On peut dire aussi que $\lambda\in\widehat Z$ est dans $Z^*$ ssi il satisfait les deux conditions :
\[
\left\{\begin{aligned}
\mathrm a)\quad &\forall(s,t)\in S,\qquad\{s,t\}\in A_0\Longleftrightarrow\{\lambda^{-1}s,\lambda t\}\in A_0\\
&[\,\text{d'où un } A_\lambda\subset\unsure{\mathfrak P|\mathcal P|P}_2(S),\ \text{défini de façon unique}\\
&\text{par la condition que }\{s,t\}\in A_0\Longleftrightarrow\{s,\lambda t\}\in A_\lambda\,]\\
\mathrm b)\quad &\text{le graphe }\Gamma_\lambda=(S,A_\lambda)\text{ est isomorphe}\\
&\text{au graphe }\Gamma_0=(S,A_0)
\end{aligned}\right.\tag*{(8)}
\]
\margin{Cela implique que $\Gamma_\lambda$ est associé à une structure \unsure{triangulée|triangul.|triangulaire} $F_\lambda$ sur $S$.}

On suppose

(9)\qquad (To 3)\qquad L'ensemble $Z^*\subset\widehat Z\subset\deleted{\unsure{G_0}}\mathfrak S_S$ est formé de \deleted{\unsure{t.tes|les}} permutations qui commutent deux à deux, \deleted{\unsure{ou|on}}

On voit (en distinguant le cas où tt sommet est déterminé par son étoile, et celui où il n'en est rien i.e. où $P_0$ est une double pyramide) que la condition (To 3)

% archive sheet number: 130
\scanpage{132}\mspage{67}

équivaut à la condition

(10)\qquad (To 3$'$)\qquad $Z^{\unsure{*|\sharp}}$ est un sous-groupe \deleted{\unsure{de $\mathfrak S_S$}}\inserted{$Z\subset\unsure{\mathfrak S|\mathcal S}_S$},

lequel sera donc nécessairement commutatif.

% slanted
\margin{\uline{NB} Il faudrait réexpliciter l'opération ``astucieuse'' de $Z$ sur $\Pi$ (p. 57, (35)), qui \uuline{n'est pas} le transport des \unsure{structure|structures}, et qu'il vaudrait donc mieux ne pas noter $i\mapsto\lambda i$, mais \unsure{avec notation|une notation} différente, genre $T_\lambda\cdot i$.}

On introduit maintenant
\[
\begin{aligned}
\Pi&\subset\unsure{\mathfrak P|\mathcal P|P}(\mathfrak P_2(S))\\[-2pt]
\rotatebox{90}{$=$}&\\[-2pt]
\{A_\lambda\mid\lambda\in Z\}&
\end{aligned}\tag*{(11)}
\]
\margin{Les groupes d'autom des $\Gamma_\lambda=(S,A_\lambda)$ sont tous égaux à $G_0$, et $Z$ opère sur $\Pi$ de façon transitive. $G_0$, $\Pi$ et $\Pi'$ ne changent pas si on remplace $(S,F_0)$ par $(S,F_i)$ $(i\in\Pi)$.}

\margin{\deleted{On avait \unsure{supposé}}\inserted{\unsure{étudié}} Pour la suite du \unsure{G|§}, on suppose $\Pi$ non réduit à un seul élément (sinon $G=G_0$, tout est trivial). Dans le cas de \unsure{(sous groupes)} la double pyramide, base $n$-gône, cela signifie que $n\equiv0\ (4)$ et $n\geqslant8$, \deleted{\unsure{auquel}}\inserted{en tous cas} dans le cas \unsure{bipyramidal|bipyramidale}, le \unsure{noyau|stabilisateur} de l'op. de $Z^*$ sur $\Pi$ est $\{1,\unsure{\underline{\circ}|\underline{-}|\varepsilon}\}$, \deleted{\illegible{1}}}

\[
\begin{aligned}
G&\subset\mathfrak S_S\\[-2pt]
\rotatebox{90}{$=$}&\\[-2pt]
\bigl\{f\in\mathfrak S_S\ \big|\ & f(A_0)\in\Pi\quad\deleted{\text{\unsure{ou encore|on}}}\ \inserted{\text{i.e. }\exists\lambda\in Z'\text{ t.q. }f(A_0)=A_\lambda}\\
&\text{i.e. } f(\Pi)=\Pi\bigr\}.
\end{aligned}\tag*{(12)}
\]

On \deleted{\unsure{fera|ferait}} note \deleted{que \unsure{les $F$|$A_i,F_i$} pour tout $\lambda\in Z$,}
\[
\deleted{\unsure{(13)}\quad\mathrm{To}\,4)}\qquad A_i\in\mathfrak P(\mathfrak P_2(S)),\quad F_i\in\mathfrak P(\mathfrak P_3(S)).\tag*{(13)}
\]
les ens. d'arêtes, resp. des faces, associés à un $i\in\Pi$ $(i=\lambda i_0,\deleted{\unsure{\forall}}\lambda\in Z)$. \inserted{il vaut mieux noter $T_\lambda i_0$} On \deleted{suppose que}\inserted{note que} $G$ ne dépend que de $\Pi\subset\mathfrak P(\mathfrak P_2(S))$ ou de l'ens. correspondant $\Pi'\subset\mathfrak P(\mathfrak P_3(S))$\unsure{;|.} On a

\diagram{\begin{tikzcd}[ampersand replacement=\&,row sep=small,column sep=small]
(14)\qquad\mathrm{Cent}(G)\subset Z\subset \& G \arrow[d,dash] \& \deleted{\text{et on}}\\
\deleted{\text{et on suppose}} \& G_0 \& {,}
\end{tikzcd}}
avec $G_0,Z$ distingués dans $G$, et on a
\[
G/G_0\lhook\joinrel\longrightarrow\mathrm{Aff}_Z(\Pi)\tag*{(15)}
\]
\margin{$[$\deleted{\unsure{Pour que}} Dans le cas bipyramidal. \unsure{p|P} avec $n\equiv0\ (4)$, $n\geqslant8$, il faut remplacer $Z$ par $Z/\{1,\underline{\circ}\}]$}
$G/G_0$ opère \uline{transitivement} sur $\Pi$

On suppose

% archive sheet number: 131
\scanpage{133}

(16)\qquad (To 4)\qquad $Z$ est central dans $G$ i.e. $Z=\mathrm{Cent}(G)$, i.e l'opération de $G/G_0$ sur $\Pi$ se fait \deleted{\unsure{(sur)}} via un hom canonique
\[
G/G_0\longrightarrow Z\qquad [\,\text{resp. }Z/\{1,\unsure{\underline{\circ}|\underline{-}|\varepsilon}\}\ \text{dans le cas bipyramidal}\,]
\]

% slanted
\margin{\uline{NB} Si cet axiome n'est pas vérifié on peut néanmoins introduire $G'=\mathrm{Ker}(G\longrightarrow\mathrm{Aut}(Z))=\mathrm{Centr}_G(Z)$, puis $Z'=\mathrm{Im}(G'\longrightarrow Z)\supset Z^2$, puis le diagramme canonique
\diagram{\begin{tikzcd}[ampersand replacement=\&,column sep=tiny,row sep=tiny]
1\arrow[r] \& G_0\arrow[r] \& G'\arrow[r]\arrow[d,dash] \& \underset{\subset Z}{Z'}\arrow[r] \& 1\\
\& \& Z\arrow[ur,swap,"\lambda\mapsto\lambda^2"] \& \&
\end{tikzcd}}
et définir aussi $\Pi'$ comme $\Pi'=\{A_\lambda\mid\lambda\in\Pi'\}$.}

de sorte qu'on a une mirifique suite exacte

\diagram{\begin{tikzcd}[ampersand replacement=\&,column sep=large]
(17)\qquad 1\arrow[r] \& G_0\arrow[r] \& G\arrow[r,"\chi'"] \& Z\arrow[r] \& 1\,,\\
\& \& Z=\mathrm{Cent}(G)\arrow[u,hook]\arrow[ur,swap,"\lambda\mapsto\lambda^2"] \& \&
\end{tikzcd}}
% slanted
\margin{dans le cas bipyramidal, il faut remplacer le $Z$ dans la suite exacte par $Z/\{1,\underline{\circ}\}$ \deleted{\unsure{mais cette suite est exacte}}}

qui est donc associée canoniquement à $(S,F_0)$ satisfaisant (To 1) à (To 4), ou encore à l'ens. $\Pi$ de structures triangulées sur $S$ défini par \deleted{$\unsure{P_0=|P_0}$} $(S,F_0)$.

Appelons les structures triangulées satisfaisant les axiomes (To 1) à (To 4) ``précongruentielles\inserted{ symplectiques}''. On introduit sur l'ens des \deleted{\unsure{relations|structures}}\inserted{structures} triangulées précongruentielles sur l'ens $S$ \unsure{une|la} relation d'équivalence

\[
\left\{\begin{aligned}
&(S,F_0)\ \text{équivalent à }(S,F_1),\quad\deleted{\text{s'il existe}}\deleted{\text{\unsure{ces str. sont}}}\inserted{\text{s'il existe}}\\
&\lambda\in Z_{F_0}\quad\text{tel que } F_1=\lambda(F_0),\quad\text{i.e. }\deleted{\text{s'il exist}}\inserted{\text{si}}\\
&\deleted{\lambda\in\unsure{\mathfrak S|\mathcal S}}\quad\mathrm a)\ (S,F_1)\ \text{isom. à }(S,F_0)\\
&\hphantom{\lambda\in\mathfrak S\quad}\mathrm b)\ \exists\lambda\in\mathfrak S_S,\ \text{telle que }\lambda\text{ commute à}\\
&\hphantom{\lambda\in\mathfrak S\quad\mathrm b)\ }G_0=\mathrm{Aut}(S,F_0),\quad\text{et que l'on ait}\\
&\hphantom{\lambda\in\mathfrak S\quad\mathrm b)\ }\{s,t\}\in A_{F_0}\Longleftrightarrow\{s,\lambda t\}\in A_{F_1}.
\end{aligned}\right.\tag*{(18)}
\]
% slanted
\margin{on pourrait élargir la définition au cas des structures ne satisfaisant pas (To 4), cf annotation marginale plus haut\ldots}

\deleted{Soit} Une \uline{structure triangulée précongruentielle}

% archive sheet number: 132
\scanpage{134}\mspage{68}

sur l'ens $S$ sera par définition la donnée d'une classe d'équivalence\inserted{ $\Pi$}, au sens de (18) de structures tr. sympl précongruentielles.

% slanted
\margin{$\Pi\hookrightarrow\mathfrak P(\mathfrak P_3(S))$ ou $\Pi\hookrightarrow\mathfrak P(\mathfrak P_2(S)$ \unsure{au|etc.} choix\ldots}

\unsure{A'|À} une telle structure est donc associée une suite exacte canonique (17), où $G_0$ est le groupe d'autom. commun des structures $(S,F_i)_{i\in\Pi}$, et où $Z$ est le groupe $Z_{F_i}$, commun pour tous les $i\in\Pi$, $\Pi$ étant de façon canonique un torseur sous $Z$.

\uline{NB} À vrai dire, cette axiomatique n'est pas encore réduite vraiment à l'os, en ce sens que les structures triangulées n'y jouent qu'un rôle logique assez accessoire -- les réflexions aboutissant aux axiomes To 3, To 4 (To 1 et To 2 $\unsure{y|{}}$ tombent) pouvant se développer le plus naturellement dans le context des \uline{graphes} plutôt que dans celui des structures triangulées qu'on peut associer à cer\uline{tains} graphe.

Pour en revenir quand $=$ aux structures triangulées, on voit que la structure \deleted{\unsure{p|P}}triangulée précongruentielle
\[
\Pi\in\mathfrak P(\mathfrak P(\mathfrak P_3(S)))\tag*{(19)}
\]

% archive sheet number: 133
\scanpage{135}

peut aussi se décrire entièrement à l'aide de
\[
\begin{aligned}
F=F_\Pi=\bigcup_{i\in\Pi}F_i=\bigl\{f\in\mathfrak P_3(S)\ \big|\ &\exists\,i\in\Pi\ \text{tel que}\\
& f\text{ soit une face}\\
&\text{pour la structure}\\
&\text{polyédrale }P_i\bigr\},
\end{aligned}\tag*{(20)}
\]
on se propose d'expliciter les conditions sur $F$, pour que $(S,F)$ provienne d'une ``structure triangulée précongruentielle'', qui sera alors néc. \uline{unique}, les $F_i$ s'interprétant alors comme les ens. des sommets des composantes connexes du graphe
\[
\Gamma_F=\{F,A,R\}\tag*{(21)}
\]
% slanted
\margin{$R$ la Relation d'inclusion $\supset$}
où on définit $A\subset\mathfrak P_2(S)$, \inserted{en termes de $F$,} comme
\[
A=\{a\in\mathfrak P_2(S)\mid\exists f\in F,\ f\supset a\}\tag*{(22)}
\]
(en fait, on aura $A=\bigcup_{i\in\Pi}A_i$) :
\[
\Pi\simeq\Pi_0(\Gamma_F).\tag*{(23)}
\]

Les axiomes sont

\deleted{(T 0)\qquad $\bigcup_{f\in F}f=S$ \unsure{(pour mémoire} -- on veut une structure triangulée !)}

(T 1)\qquad $\deleted{\unsure{\exists}}\forall\,a\in A$, existe exactement \uline{deux} $f\in F$ qui contiennent $a$,

ce qui permet de considérer $\Gamma_F$ comme un graphe. Prenant les composantes connexes

% archive sheet number: 134
\scanpage{136}\mspage{69}

de ce dernier, on trouve une décomposition
\[
F=\coprod_{i\in\Pi}F_i,\qquad A=\coprod_{i\in\Pi}A_i\tag*{(\deleted{$\unsure{T}$}24)}
\]
les $F_i,A_i$ non vides et mutuellement disjoints.

On exige alors (cf p. 53)

(25)\qquad (T 3)\qquad $\forall i\in\Pi$, $(S,F_i)$ est une structure \uline{triangulée} \uline{stricte}\inserted{ $P_i$}, i.e. $\bigcup_{f\in F_i}f=S$, et on a connexité du contour local en $s$ défini par $P_i$,

cf T3 a) b) p. $53'$ pour une explicitation.

Les autres axiomes peuvent s'exprimer, en disant que la famille $(F_i)$ des structures triangulées sur $S$ forme exactement une classe d'équivalence de structures triangulées symplectiques précongruentielles.

% slanted
\margin{En \unsure{travaillant|traduisant|dimension} avec les structures triangulées $(S,F)$ de $=$, on trouve que (25 bis) $\left\{\begin{aligned}&G\text{ est simplement transitif}\\&\text{sur l'ens. des ``\uline{repères}''}\\&\text{de }(S,F)\end{aligned}\right.$ donc il est transitif sur $\vec F,F,\vec A,A,S$, le stabilisateur d'un él. $\vec f$ resp $f$ resp. $\vec a$ resp $a$ resp $s$ est isomorphe à $\mathbb Z/3\mathbb Z$ (can.), $\unsure{\mathbb D|D}_3$, $\mathbb Z/2\mathbb Z$, $\mathbb D_{\unsure{2|{}}}^{2}$, et est \unsure{isom|isom.} (pour $s$) à un produit $\frac12$ direct}

Il faut en arriver enfin, progressivement, à des conditions plus strictes encore, pour en arriver enfin aux \inserted{structures} \deleted{qu'on aura envie d'appeler} ``congruentielles'', isomorphes à celles déduites des modules libres de rang 2 sur des $\mathbb Z/n\mathbb Z$. Pour l'instant, l'attention principale \deleted{\unsure{est}} sera sur les phénomènes liés au cas d'un $\Pi$ non réduit à un point, i.e. d'un $Z$ qui n'est pas réduit à $\{1\}$ (ou à $\{1,\unsure{\underline{\circ}|\underline{-}|\varepsilon}\}$, dans le cas

% archive sheet number: 135
\scanpage{137}

bipyramidal). Je voudrais \inserted{en} arriver à définir un isomorphisme
\[
\underset{(\text{resp }Z/\{1,\unsure{\underline{\circ}|\underline{-}|\varepsilon}\})}{Z}\xrightarrow{\sim}(\mathbb Z/n\mathbb Z)^{*}/\{\pm1\}\tag*{(26)}
\]
pour $n$ convenable.

On va donc revenir, pour fixer les idées, à la donnée d'une structure \unsure{triangulée|striangulée}-symplectique précongruentielle. On supposera

(27)\qquad (To 5)\qquad Les sommets de $S$ sont tous de $=$ ordre fini, soit $n\in\mathbb N$ $(n\geqslant3)$.

\unsure{Soits|Soit|Soient} maintenant $\lambda\in Z$, $s\in S$, on a donc par définition de $A_\lambda$
\[
\mathrm{Et}_{A_0}(s)=\mathrm{Et}_{A_\lambda}(\lambda s)\tag*{(28)}
\]
donc l'ens \inserted{$\unsure{\mathcal E|\mathscr E|E}_0$} des étoiles\inserted{ pour $A_0$,} des \deleted{points} sommets $s\in S$ est le $=$ que l'ens $\mathcal E_\lambda$ pour $A_\lambda$ -- simplement les ``centres'' de ces étoiles changent, quand on change de structure $A_i$ $(i\in\Pi)$ -- et bien entendu aussi, la structure polygonale.

Notons qu'une structure triangulée $(S,F_0)$ satisfaisant \deleted{(\unsure{To 1})} (To 2) est déterminée par la connaissance de l'ens $\mathcal E^{\unsure{s|{}}}$ des étoiles qu'elle définit, ici en vertu de (27)
\[
\mathcal E_0\subset\mathfrak P_n(S),\qquad \deleted{(\unsure{\mathrm{NB}}\quad E\in\mathcal E_0\Rightarrow\mathrm{Card}\,E=n}\tag*{(29)}
\]
et avec la donnée en plus, \deleted{pour tte \unsure{étoile}} \inserted{de}

% archive sheet number: 136
\scanpage{138}\mspage{70}

En effet, \deleted{une telle structure $n$-polygonale est définie en termes} \inserted{prenons d'abord des structures} $n$-polygônales orientées, soit $\overrightarrow{\mathrm{Pol}}(E)$, alors on a

(33)\qquad Soit $E$ de cardinal $n$, alors $(\mathbb Z/n\mathbb Z)^*$ opère canoniquement sur $\overrightarrow{\mathrm{Pol}}(E)$.

En effet, un élément de $\unsure{\overline{\mathrm{Pol}}|\overrightarrow{\mathrm{Pol}}}(E)$ s'interprète comme
\[
\begin{aligned}
&\text{opération de }\mathbb Z/n\mathbb Z=\Lambda_n\text{ sur }E\\
&\text{faisant de }E\text{ un }\Lambda_n\text{-torseur}\\
&\text{(i.e. une ``droite affine sur }\Lambda_n\text{'')}
\end{aligned}
\]
donc $\mathrm{Aut}(\Lambda_n)$ opère sur l'ens de ces structures, par composition. Si on identifie un $\omega\in\overrightarrow{\mathrm{Pol}}(E)$ comme une permutation circulaire, on aura
\[
\begin{aligned}
&\deleted{\unsure{\omega}}\ \underbrace{T_{\nu}\omega}=\omega^\nu\\
&\qquad\text{effet sur }\omega\in\overrightarrow{\mathrm{Pol}}(E)\subset\unsure{\mathfrak S|\mathcal S}_E\\
&\qquad\text{de }T_\nu\text{ associé à }\nu\in(\mathbb Z/n\mathbb Z)^*
\end{aligned}\tag*{(34)}
\]
On a donc
\[
\left\{\begin{aligned}
&T_{-\nu}\omega=\omega^{-\nu}=(\omega^\nu)^{-1}=(T_\nu\omega)^{-1}\\
&\qquad\text{structure polygonale orientée}\\
&\qquad\text{opposée à }\deleted{\unsure{\omega}}T_\nu\omega\\
&T_\nu\omega^{-1}=(T_\nu\omega)^{-1}\qquad\text{i.e. }T_\nu\text{ commute}\\
&\qquad\qquad\deleted{\unsure{\text{à}}}\enspace\omega\longmapsto\omega^{-1}\\
&\qquad\qquad\deleted{\text{\unsure{passage}}}\quad\text{passage d'une str.}\\
&\qquad\qquad\text{de polyg. orient. à \unsure{l'opposée|l'opposé}}
\end{aligned}\right.\tag*{(35)}
\]

% archive sheet number: 137
\scanpage{139}

(30)\qquad $\forall E\in\mathcal E_0$, \deleted{\unsure{une}}la donnée de la structure de $n$-polygone combinatoire, ayant $E$ comme ens. de ses sommets, i.e. \inserted{de} $A_E\subset\mathfrak P_2(E)$

puisqu'on récupère $A_0$ à partir de (29) (30) par
\[
A_0=\{a\in\mathfrak P_2(S)\mid\exists E\in\mathcal E_0\ \text{t.q. }a\in A_E\}.\tag*{(31)}
\]
À vrai dire, \deleted{pour} dans le cas où l'axiome (To 1) est satisfait, il existe exactement \uline{deux} $s\in S$ tels que $\mathrm{Et}_{\deleted{\unsure{\lambda}}A_0}(s)$ contienne $a$ (donc admette $a$ comme arête) et sauf dans le cas \unsure{octaèdral|octaédral}, ces étoiles associées sont distinctes, donc sauf dans le cas de l'octaèdre, il y aura dans (31) exactement deux $E\in\mathcal E_0$ telles que $E\supset a$, i.e $a\in A_E$. (Dans le cas de l'octaèdre, par contre, le $E$ en question est unique.).

Pour déterminer \deleted{\unsure{une}} \unsure{structure|structures} $A_\lambda$ $(\lambda\in Z)$, comme \inserted{l'ens $\mathcal E_\lambda$ de} ses étoiles est égal à $\mathcal E_0$, il suffit de savoir décrire, \deleted{$\unsure{E}$} pour $\forall E\in\mathcal E_0$, la \deleted{\unsure{struc}}\inserted{structure} $n$-polygonale $A_{E,\lambda}\in\mathfrak P_2(E)$ induite sur $E$ par $\Gamma_\lambda=(S,A_\lambda)$. Notons maintenant le fait mirifique

(32)\qquad Soit $E$ un ens. \inserted{fini,} de cardinal $n\geqslant3$. Alors le groupe $Z'_n=(\mathbb Z/n\mathbb Z)^*/\{\pm1\}$ opère canoniquement sur l'ens. \inserted{$\mathrm{Pol}(E)$} des structures polygônales sur $E$ comme ens de sommets :

% slanted
\margin{On voit de suite que cette opération est fidèle.}

% archive sheet number: 138
\scanpage{140}\mspage{71}

donc l'opération de $(\mathbb Z/n\mathbb Z)^*$ sur $\overrightarrow{\mathrm{Pol}}(E)$ passe au quotient en une opération sur
\[
\deleted{\mathrm{Pol}(E)}\qquad\mathrm{Pol}(E)=\overrightarrow{\mathrm{Pol}}(E)/(\omega\mapsto\omega^{\unsure{\nu|-1}}),
\]
enfin l'élément $-1\in(\mathbb Z/n\mathbb Z)^*$ opère trivialement sur $\mathrm{Pol}(E)$, en vertu de 35 a), d'où l'opération (32).

Pour $\nu\in(\mathbb Z/n\mathbb Z)^*/\pm1=Z'_n$, et $\alpha$ structure $n$-polygonale sur un ens $E$, on désigne par $\alpha^\nu$ la structure \deleted{\unsure{induite}}\inserted{transformée} par $\nu$. Ceci fixé, j'en viens à l'axiome massue :

% slanted
\margin{Si $A\subset\mathfrak P_2(E)$ définit la structure polygonale, on désigne par $A^{(\nu)}\subset\mathfrak P_2(E)$ sa transformée par $\nu\in Z'_n$.}

(33)\qquad (To 6)\qquad $\forall\lambda\in Z$, \deleted{$\unsure{s\in S},\ \unsure{\varepsilon}$} \inserted{$\exists\nu=\nu(\lambda)\in\underbrace{(\mathbb Z/n\mathbb Z)^*/\pm1}_{Z'_n}$} (évidemment unique) tel que pour tout $s\in S$, la structure $n$-polygonale \deleted{induite} $\alpha_{\lambda.i_0,\lambda.\unsure{s|b}}$ induite par $A_{\lambda.i_0}=A_\lambda$\inserted{ (plutôt : $T_\lambda i_0$)} sur $\mathrm{Et}_{A_0}(s)=\mathrm{Et}_{A_\lambda}(\lambda s)$, soit égale à $(\alpha_{i_0,s})^\nu$.

Cela détermine donc la structure $A_\lambda$ en termes de $\nu=\nu(\lambda)\in Z'_n$, \inserted{notons-la $A_0^{(\nu)}$,} donc dans le cas où $P_0$ n'est pas l'octaèdre (\deleted{\unsure{au cas}}\inserted{dans lequel} de toutes façons $Z=\{1,\unsure{\underline{\circ}|\underline{-}|\varepsilon}\}$ opère trivialement sur $\Pi$, réduit à un pt !), cela détermine aussi $\lambda$ en termes de $\nu$ :

% archive sheet number: 139
\scanpage{141}

(34)\qquad $\lambda\deleted{\unsure{(s)}}s$ est l'unique élément de $S$, dont l'étoile pour $A_\lambda=A_0^{\nu}$ est égale à $\mathrm{Et}_0(s)$.

$[$\uline{NB} Dans le cas de l'octaèdre, rappelons que $Z=\{1,\unsure{-|\underline{\circ}|\varepsilon}\}$, et que pour $s\in S$, $\unsure{\underline{\circ}|\underline{-}|\varepsilon}(s)$ peut se caractériser comme l'unique $s'\ne s$ \deleted{qui} \unsure{dont|qui} a même étoile que $s$. $]$ Ainsi on trouve un \uline{homomorphisme} \uline{injectif}
\[
Z\hookrightarrow\Lambda'_n\tag*{(35)}
\]
\margin{(cas différent du cas octaédral \deleted{\unsure{y|qui}} où on a $\Lambda'_n=\{1\}$, $Z=\{1,\underline{\circ}\}$)}
dont l'image est formée d'éléments $\nu$ satisfaisant la condition suivante

(36)\qquad L'ensemble \deleted{$\unsure{A'_0|A_0^{(\nu)}}=\{a\in\unsure{\mathfrak P|\mathcal P|P}_2(S)\mid\exists E\in\unsure{\mathcal E|\mathscr E|E}_0\}$} \deleted{$\unsure{\coprod_{E\in\mathcal E_0} A_E\text{ t.q. }a\in{}}$}
\[
A_0^{(\nu)}=\coprod_{E\in\mathcal E_0}A_E^{(\nu)}
\]
est l'ensemble des arêtes pour une structure de 2-polyèdre sur $S$ $[$ayant mêmes ``étoiles que $P_0$, et en fait, plus précisément, \uline{équivalente} à celle de $P_0$ au sens de (18)$]$.

Mais le simple fait déjà, \deleted{que} pour une structure de 2-polyèdre triangulé $P_0=(S,A_0,F_0)$ satisfaisant $[$To 1 et$]$ To 2,\inserted{ (To 5),} que l'ens $A_0^{(\nu)}\subset\mathfrak P_2(S)$ définisse sur $S$ une structure de graphe provenant, elle aussi, d'une telle structure triangulée -- et ceci \uline{pour tout $\nu\in\deleted{\unsure{Z'_n}}\Lambda'_n$} -- me paraît déjà une propriété remarquable et très forte, qui mérite d'être examinée

% slanted
\margin{cf To 7 plus bas !}

% archive sheet number: 140
\scanpage{142}\mspage{72}

pour son propre mérite. La question bien sûr ne se pose que $Z'_n\ne\{1\}$, i.e. si $n\notin\{2,3,4,6\}$, le premier cas est donc celui où
\[
n=5,\qquad \Lambda'_{n}\simeq\deleted{\unsure{\Lambda}}\mathbb F_5^*/\pm1\simeq\mathbb Z/2\mathbb Z,\tag*{(37)}
\]
considérons une structure de 2-polyèdre connexe triangulé ($=$ faces d'ordre 3), dont les sommets sont \inserted{tous} d'ordre 5. On sait qu'il n'y a que l'icosaèdre (structure orientable à 12 sommets) et l'icosaèdre gauche (structure non orientée). L'opération envisagée ici (pour $\unsure{\nu|v}$ l'élément non trivial de $\Lambda'_5$) est celle qui fait passer d'un icosaèdre à l'icosaèdre associé -- l'ens $\Pi$ (dans le cas orienté) étant la structure de \uline{bi-icosaèdre}.

L'icosaèdre peut être décrit à l'aide d'un vectoriel $M$ de rang 2 sur $\mathbb F_5$, comme $(S_M,F_M)$, et satisfait donc à toute l'axiomatique du présent §. Mais l'icosaèdre gauche ? Notons que le commutant de $G_0$ dans ce cas est réduit à l'identité, donc $Z=1$, donc on ne s'attendra pas à ce que $Z\longrightarrow\deleted{\unsure{Z}}\Lambda_n^{'\!*}$ soit surjectif -- pourtant l'opération $A_0\longrightarrow A_0^{(\nu)}$ a un sens dans ce cas là ! Cela

% archive sheet number: 141
\scanpage{143}

montre bien que dans certains cas, la validité de l'opération $A_0\longmapsto A_0^{\unsure{\nu|v|\vee}}$ déborde de celle de l'axiomatique $\mathrm{To}\,1:\mathrm{To}\,5$ développée jusqu'à présent.

Le cas suivant serait
\[
 n=7,\qquad Z'_n\simeq\mathbb Z/3\mathbb Z
 \qquad \deleted{\unsure{\simeq\mathbb Z|=}}\quad
 ((\mathbb Z/7\mathbb Z)^\times\simeq\mathbb Z/6\mathbb Z)
\]
engendré par n'importe quel
\[
 \unsure{\zeta|\nu|\xi}\in\mathbb F_7^\times\setminus\deleted{\unsure{\{\pm\}|\{1\}}}\{1,-1\}
 =\mathbb F_7^{\deleted{\unsure{\times|*}}}\setminus\{0,1,-1\}
\]
p.ex. $\nu=2\bmod 7$, dont l'inverse est $\nu'=4\bmod 7$,

est-il vrai que pour tout 2-polyèdre \deleted{\unsure{rég|rég.}} triangulé \inserted{$(S,A_0,F_0)$} dont les sommets sont \inserted{tous} d'ordre 7, \unsure{l'|le} \deleted{\unsure{g|gr}} gr. pour $(S,A_0^{(2)}\subset\mathfrak P_2(S))$ correspond encore à un tel 2-polyèdre ? Même question pour $n=8$, avec $Z'_8\simeq\mathbb Z/2\mathbb Z$, ou \inserted{$n=10$ avec} $Z'_{10}\simeq\mathbb Z/2\mathbb Z$ (pour $n=9$ on a $Z'_9\simeq\mathbb Z/4\mathbb Z$ -- en fait pour
\[
\left\{\begin{aligned}
&n\notin\{\underbrace{2,3,4,6}_{Z'_n=1},
 \underbrace{5,8,10\unsure{\}|{}},12}_{Z'_n\simeq\mathbb Z/2\mathbb Z})
 \quad\text{i.e.}\quad n\geq\unsure{18|13}\ \text{ou}\ n\in\{
 \underbrace{7}_{Z'_n\simeq\mathbb Z/3\mathbb Z},
 \underbrace{9}_{Z'_n\in\mathbb Z/4\mathbb Z},11\},\ \text{on a}\\
&Z'_n\ \text{d'ordre}\ \geq3\ \text{-- ordre 3 pour}\ n=7\ \text{et }14,\ \text{ordre }4\\
&\text{pour }n=9\ \text{et }18\ \inserted{\text{et pour }n=16,20}\ \text{-- ordre 5 pour }n=11\ \text{et }22,\ \text{-- ordre 6 pour}\\
&n=13\ \text{et }26\ \text{-- }\inserted{\text{(ordre 12 pour }n=15\ \text{ou }30\text{)}}\
  \deleted{\text{\unsure{ordre|ordre 8}}}\ \text{ordre}\\
&8\ \text{pour }n=17\ \text{ou }34,\ \deleted{\text{\unsure{10 pour|ordre}}}\
 \inserted{\text{9 pour}}\ 19\ \text{et }38,
\end{aligned}\right.
\]

Mais poursuivons sur la lancée \unsure{des|de} l'axiomatique des $\mathrm{To}\,1\ldots$, on trouve
\[
 \text{(To 7)}\quad \text{L'homomorphisme}\quad Z\hookrightarrow\Lambda'_n\quad\text{est surjectif.}\tag*{(38)}
\]
i.e. un isomorphisme.

% archive sheet number: 142
\scanpage{144}\mspage{73}

\[
 \text{(To 8)}\qquad\text{Le 2-polyèdre $P_0$ est régulier, et orientable}\tag*{(39)}
\]
\margin{Il vaut mieux séparer en deux axiomes distincts (39) -- ils ne jouent pas le $=$ rôle. De plus, il semble que dans la plupart des arguments, on n'ait besoin que du fait que $G_0^+$ soit transitif sur l'ens \unsure{$\swarrow$|.} \deleted{\unsure{$S$|$F$}} \unsure{??|?}}% slanted

ce qui fait qu'on doit pouvoir traduire notre axiomatique en termes de théorie des groupes, \unsure{d'abord|à savoir|savoir} en termes des propriétés d'un groupe $G_0^+$ qui est engendré par des générateurs $\rho_0,\rho_1,\rho_\infty$ satisfaisant
\[
 \boxed{\rho_0^n=\rho_1^2=\rho_\infty^3=1}\tag*{(40)}
\]
et où $\rho_0$ est d'ordre $n$ exactement, sous-groupe \deleted{\unsure{de|d'un}} d'indice 2 d'un groupe \inserted{$G_0$} ayant des générateurs $\tau_0,\tau_1,\tau_\infty$ reliés entre eux et aux précédents par
\[
 \left\{\begin{aligned}
 &\tau_0^2=\tau_1^2=\tau_\infty^2=1\\
 &\rho_0=\tau_\infty\tau_1,\quad \rho_1=\tau_0\tau_\infty,\quad
 \rho_\infty=\tau_1\tau_0 .
 \end{aligned}\right. \tag*{(41)}
\]

Je vais d'abord \deleted{introduire un \unsure{homo|bon} isomorphisme} \inserted{introduire le $\Lambda_n$-torseur}
\[
 \vec\Pi=\deleted{\unsure{\coprod|=}}\coprod_{i\in\Pi}\mu(i)\qquad
 \text{où }\mu(i)=\text{ens. des deux orientations de }P_i=(S,A_i,F_i)\tag*{(42)}
\]
on fait opérer dessus \deleted{$(\mathbb Z/n\mathbb Z)^*\simeq\unsure{\Lambda|\mathbb Z}$}
\[
 Z_n=\Lambda_n^{*}=(\mathbb Z/n\mathbb Z)^*\tag*{(43)}
\]
en notant que pour tout \deleted{$s\in S$ et} \inserted{$i\in\Pi$,} $\lambda\in Z'\simeq Z'_n$, \inserted{$s\in S$} on a sur $\underbrace{\operatorname{Et}_i(s)=\operatorname{Et}_{\lambda.i}(\lambda s)}_{E}$, deux structures polygônales \deleted{\unsure{$\alpha$|$a_i$}} $\alpha_i(s)$, $\alpha_{\lambda.i}(s)$, reliées par
\[
 \alpha_{\lambda.i}(s)=\alpha_i(s)^\lambda\tag*{(44)}
\]

% archive sheet number: 143
